\documentclass[12pt,a4paper]{article}

\usepackage[utf8]{inputenc}
\usepackage[T1]{fontenc}
\usepackage[polish]{babel}
\usepackage{geometry}
\usepackage{amsmath}
\usepackage{booktabs}
\usepackage{float}
\usepackage{array}
\usepackage{enumitem}

\geometry{margin=2.5cm}

\title{Sprawozdanie laboratoryjne nr 3\\
\large Inteligencja obliczeniowa i jej zastosowania\\
\large Wersja A -- metoda AHP}
\author{Jakub Jaśków 268416, Justyna Ziemichód 268431}
\date{23 marca 2026}

\begin{document}

\maketitle

\section{Cel ćwiczenia}

Celem ćwiczenia było zdefiniowanie wielokryterialnego problemu wyboru produktu oraz rozwiązanie go metodą AHP (\textit{Analytic Hierarchy Process}). Zgodnie z poleceniem należało:
\begin{itemize}
    \item określić kryteria najwyższego poziomu,
    \item podać subkryteria przynajmniej dla jednego kryterium,
    \item wskazać alternatywy,
    \item wykonać obliczenia dla kryteriów najwyższego poziomu i co najmniej jednego subkryterium,
    \item sprawdzić współczynnik spójności \(CR\).
\end{itemize}

\section{Opis problemu decyzyjnego}

Rozważanym problemem jest wybór tableta spośród czterech dostępnych alternatyw. Paramtetry poszczególnych tabletów zostały okreśone na podstawie prawdziwych, rzeczywistych urządzeń. W dalszej części sprawozdania, dla uproszczenia zapisu obliczeń, ich nazwy będą zastępowane oznaczeniami \textbf{Tablet 1}, \textbf{Tablet 2}, \textbf{Tablet 3} oraz \textbf{Tablet 4}.

\subsection*{Alternatywy}
\begin{itemize}
    \item A1 -- Tablet 1 -- Lenovo Tab Pro 12{,}7'' 8GB/256GB
    \item A2 -- Tablet 2 -- REDMAGIC Astra 16GB/512GB
    \item A3 -- Tablet 3 -- Samsung Galaxy Tab S11 Ultra 16GB/1TB
    \item A4 -- Tablet 4 -- OneXPlayer X1 Air 32GB/1TB
\end{itemize}

\subsection*{Specyfikacja analizowanych tabletów}

\begin{table}[H]
\centering
\caption{Specyfikacja analizowanych tabletów}
\begin{tabular}{lcccc}
\toprule
\textbf{Tablet} & \textbf{Cena} & \textbf{Pamięć RAM} & \textbf{Pojemność dysku} & \textbf{Czas pracy na baterii} \\
\midrule
Tablet 1 & 2500 zł & 8 GB  & 256 GB  & 6 h \\
Tablet 2 & 3200 zł & 16 GB & 512 GB  & 8 h \\
Tablet 3 & 4200 zł & 16 GB & 1000 GB & 10 h \\
Tablet 4 & 5200 zł & 32 GB & 1000 GB & 12 h \\
\bottomrule
\end{tabular}
\end{table}

\subsection*{Kryteria najwyższego poziomu}

Przyjęto cztery kryteria:
\begin{itemize}
    \item \(K_1\) -- Cena,
    \item \(K_2\) -- Pamięć RAM,
    \item \(K_3\) -- Pojemność dysku,
    \item \(K_4\) -- Czas pracy na baterii.
\end{itemize}

\subsection*{Subkryteria}

Dla zachowania pełnej hierarchii decyzyjnej przyjęto następujące subkryteria:

\begin{itemize}
    \item dla kryterium \(K_1\) -- Cena:
    \[
    S_{1}=\{2500\text{ zł}, 3200\text{ zł}, 4200\text{ zł}, 5200\text{ zł}\}
    \]
    \item dla kryterium \(K_2\) -- Pamięć RAM:
    \[
    S_{2}=\{8\text{ GB}, 16\text{ GB}, 32\text{ GB}\}
    \]
    \item dla kryterium \(K_3\) -- Pojemność dysku:
    \[
    S_{3}=\{256\text{ GB}, 512\text{ GB}, 1000\text{ GB}\}
    \]
    \item dla kryterium \(K_4\) -- Czas pracy na baterii:
    \[
    S_{4}=\{6\text{ h}, 8\text{ h}, 10\text{ h}, 12\text{ h}\}
    \]
\end{itemize}

\section{Hierarchia problemu}

Hierarchia decyzyjna ma postać:
\[
\text{Cel} \rightarrow \text{Kryteria} \rightarrow \text{Subkryteria} \rightarrow \text{Alternatywy}
\]

czyli:
\[
\text{wybór tableta} \rightarrow (K_1,K_2,K_3,K_4) \rightarrow (S_1,S_2,S_3,S_4) \rightarrow (A_1,A_2,A_3,A_4)
\]

\section{Metoda AHP}

W metodzie AHP buduje się macierz porównań parami:
\[
A=[a_{ij}]
\]

Następnie wykonuje się normalizację kolumn:
\[
b_{ij}=\frac{a_{ij}}{\sum_{i=1}^{n} a_{ij}}
\]

Wagi oblicza się jako średnie arytmetyczne wierszy znormalizowanej macierzy:
\[
w_i=\frac{1}{n}\sum_{j=1}^{n} b_{ij}
\]

Spójność ocen sprawdza się następująco:
\[
V=\frac{A\cdot W}{W}
\]
\[
\lambda_{\max}=\frac{1}{n}\sum_{i=1}^{n} v_i
\]
\[
CI=\frac{\lambda_{\max}-n}{n-1}
\]
\[
CR=\frac{CI}{RI}
\]

Jeżeli:
\[
CR<0{,}1
\]
to macierz uznaje się za wystarczająco spójną.

\section{Wyznaczenie wag kryteriów najwyższego poziomu}

Przyjęto, że przy wyborze tableta najważniejsza jest \textbf{Cena}, następnie \textbf{Pamięć RAM}, potem \textbf{Pojemność dysku}, a na końcu \textbf{Czas pracy na baterii}.

Oznacza to, że:
\begin{itemize}
    \item cena jest umiarkowanie ważniejsza od pamięci RAM,
    \item cena jest ważniejsza od pojemności dysku,
    \item cena jest o wiele ważniejsza od czasu pracy na baterii,
    \item pamięć RAM jest umiarkowanie ważniejsza od pojemności dysku,
    \item pamięć RAM jest ważniejsza od czasu pracy na baterii,
    \item pojemność dysku jest umiarkowanie ważniejsza od czasu pracy na baterii.
\end{itemize}

\begin{table}[H]
\centering
\caption{Tabela porównań parami kryteriów najwyższego poziomu}
\begin{tabular}{lcccc}
\toprule
 & \textbf{Cena} & \textbf{Pamięć RAM} & \textbf{Pojemność dysku} & \textbf{Czas pracy na baterii} \\
\midrule
\textbf{Cena} & 1 & 3 & 5 & 7 \\
\textbf{Pamięć RAM} & \(1/3\) & 1 & 3 & 5 \\
\textbf{Pojemność dysku} & \(1/5\) & \(1/3\) & 1 & 3 \\
\textbf{Czas pracy na baterii} & \(1/7\) & \(1/5\) & \(1/3\) & 1 \\
\bottomrule
\end{tabular}
\end{table}

Na tej podstawie zbudowano macierz porównań kryteriów:

\[
A_K=
\begin{bmatrix}
1 & 3 & 5 & 7 \\
1/3 & 1 & 3 & 5 \\
1/5 & 1/3 & 1 & 3 \\
1/7 & 1/5 & 1/3 & 1
\end{bmatrix}
\]

Kolejność kryteriów w macierzy:
\[
(K_1,K_2,K_3,K_4)=(\text{Cena},\text{RAM},\text{Dysk},\text{Bateria})
\]


\subsection*{Krok 1. Sumy kolumn}

\[
\sum K_1 = 1+\frac13+\frac15+\frac17 = 1{,}6762
\]
\[
\sum K_2 = 3+1+\frac13+\frac15 = 4{,}5333
\]
\[
\sum K_3 = 5+3+1+\frac13 = 9{,}3333
\]
\[
\sum K_4 = 7+5+3+1 = 16
\]

\subsection*{Krok 2. Macierz znormalizowana}

\[
B_K=
\begin{bmatrix}
0{,}5966 & 0{,}6618 & 0{,}5357 & 0{,}4375 \\
0{,}1989 & 0{,}2206 & 0{,}3214 & 0{,}3125 \\
0{,}1193 & 0{,}0735 & 0{,}1071 & 0{,}1875 \\
0{,}0852 & 0{,}0441 & 0{,}0357 & 0{,}0625
\end{bmatrix}
\]

\subsection*{Krok 3. Wektor wag}

Średnie wierszy dają wektor wag:
\[
W=
\begin{bmatrix}
0{,}5579\\
0{,}2633\\
0{,}1219\\
0{,}0569
\end{bmatrix}
\]

Zatem:
\begin{itemize}
    \item Cena -- \(0{,}5579\)
    \item Pamięć RAM -- \(0{,}2633\)
    \item Pojemność dysku -- \(0{,}1219\)
    \item Czas pracy na baterii -- \(0{,}0569\)
\end{itemize}

\section{Analiza zgodności dla kryteriów najwyższego poziomu}

Najpierw obliczono:
\[
A_K \cdot W =
\begin{bmatrix}
2{,}3555\\
1{,}0994\\
0{,}4919\\
0{,}2299
\end{bmatrix}
\]

Następnie:
\[
V=
\frac{A_K \cdot W}{W}
=
\begin{bmatrix}
4{,}2222\\
4{,}1747\\
4{,}0362\\
4{,}0408
\end{bmatrix}
\]

\[
\lambda_{\max}=\frac{4{,}2222+4{,}1747+4{,}0362+4{,}0408}{4}=4{,}1185
\]

\[
CI=\frac{4{,}1185-4}{4-1}=0{,}0395
\]

Dla macierzy \(4\times 4\):
\[
RI=0{,}90
\]

\[
CR=\frac{0{,}0395}{0{,}90}=0{,}0439
\]

Ponieważ:
\[
CR=0{,}0439<0{,}1
\]
macierz porównań kryteriów najwyższego poziomu jest spójna.

\section{Wyznaczenie wag i analiza zgodności subkryteriów}

\subsection*{1. Subkryteria dla kryterium Cena}

Dla kryterium \textbf{Cena} przyjęto, że niższa cena jest korzystniejsza:
\[
2500\text{ zł} > 3200\text{ zł} > 4200\text{ zł} > 5200\text{ zł}
\]

\begin{table}[H]
\centering
\caption{Tabela subkryteriów dla kryterium Cena}
\begin{tabular}{lcccc}
\toprule
 & \textbf{2500 zł} & \textbf{3200 zł} & \textbf{4200 zł} & \textbf{5200 zł} \\
\midrule
\textbf{2500 zł} & 1 & 3 & 5 & 7 \\
\textbf{3200 zł} & 1/3 & 1 & 3 & 5 \\
\textbf{4200 zł} & 1/5 & 1/3 & 1 & 3 \\
\textbf{5200 zł} & 1/7 & 1/5 & 1/3 & 1 \\
\bottomrule
\end{tabular}
\end{table}

\[
A_{S_1}=
\begin{bmatrix}
1 & 3 & 5 & 7 \\
1/3 & 1 & 3 & 5 \\
1/5 & 1/3 & 1 & 3 \\
1/7 & 1/5 & 1/3 & 1
\end{bmatrix}
\]

\[
W=
\begin{bmatrix}
0{,}5579\\
0{,}2633\\
0{,}1219\\
0{,}0569
\end{bmatrix}
\]

\[
A_{S_1}\cdot W=
\begin{bmatrix}
2{,}3555\\
1{,}0994\\
0{,}4919\\
0{,}2299
\end{bmatrix}
\]

\[
V=
\frac{A_{S_1}\cdot W}{W}
=
\begin{bmatrix}
4{,}2222\\
4{,}1747\\
4{,}0362\\
4{,}0408
\end{bmatrix}
\]

\[
\lambda_{\max}=4{,}1185,\qquad
CI=0{,}0395,\qquad
CR=0{,}0439
\]

Zatem macierz subkryteriów ceny jest spójna.

\subsection*{2. Subkryteria dla kryterium Pamięć RAM}

Dla pamięci RAM przyjęto, że większa wartość jest korzystniejsza:
\[
32\text{ GB} > 16\text{ GB} > 8\text{ GB}
\]

\begin{table}[H]
\centering
\caption{Tabela subkryteriów dla kryterium Pamięć RAM}
\begin{tabular}{lccc}
\toprule
 & \textbf{8 GB} & \textbf{16 GB} & \textbf{32 GB} \\
\midrule
\textbf{8 GB} & 1 & 1/3 & 1/5 \\
\textbf{16 GB} & 3 & 1 & 1/3 \\
\textbf{32 GB} & 5 & 3 & 1 \\
\bottomrule
\end{tabular}
\end{table}

\[
A_{S_2}=
\begin{bmatrix}
1 & 1/3 & 1/5 \\
3 & 1 & 1/3 \\
5 & 3 & 1
\end{bmatrix}
\]

Suma kolumn:
\[
\sum S_{2,1}=9,\qquad \sum S_{2,2}=4{,}3333,\qquad \sum S_{2,3}=1{,}5333
\]

Macierz znormalizowana:
\[
B_{S_2}=
\begin{bmatrix}
0{,}1111 & 0{,}0769 & 0{,}1304 \\
0{,}3333 & 0{,}2308 & 0{,}2174 \\
0{,}5556 & 0{,}6923 & 0{,}6522
\end{bmatrix}
\]

Wektor wag:
\[
W=
\begin{bmatrix}
0{,}1061\\
0{,}2605\\
0{,}6334
\end{bmatrix}
\]

\[
A_{S_2}\cdot W=
\begin{bmatrix}
0{,}3199\\
0{,}7906\\
1{,}9453
\end{bmatrix}
\]

\[
V=
\frac{A_{S_2}\cdot W}{W}
=
\begin{bmatrix}
3{,}0151\\
3{,}0349\\
3{,}0713
\end{bmatrix}
\]

\[
\lambda_{\max}=3{,}0404
\]

\[
CI=\frac{3{,}0404-3}{2}=0{,}0202
\]

Dla macierzy \(3\times 3\):
\[
RI=0{,}58
\]

\[
CR=\frac{0{,}0202}{0{,}58}=0{,}0348
\]

Ponieważ:
\[
CR<0{,}1
\]
macierz subkryteriów pamięci RAM jest spójna.

\subsection*{3. Subkryteria dla kryterium Pojemność dysku}

Dla pojemności dysku przyjęto, że większa wartość jest korzystniejsza:
\[
1000\text{ GB} > 512\text{ GB} > 256\text{ GB}
\]

\begin{table}[H]
\centering
\caption{Tabela subkryteriów dla kryterium Pojemność dysku}
\begin{tabular}{lccc}
\toprule
 & \textbf{256 GB} & \textbf{512 GB} & \textbf{1000 GB} \\
\midrule
\textbf{256 GB} & 1 & 1/3 & 1/5 \\
\textbf{512 GB} & 3 & 1 & 1/3 \\
\textbf{1000 GB} & 5 & 3 & 1 \\
\bottomrule
\end{tabular}
\end{table}

\[
A_{S_3}=
\begin{bmatrix}
1 & 1/3 & 1/5 \\
3 & 1 & 1/3 \\
5 & 3 & 1
\end{bmatrix}
\]

Ponieważ macierz ma taką samą strukturę jak dla pamięci RAM, otrzymujemy:
\[
W=
\begin{bmatrix}
0{,}1061\\
0{,}2605\\
0{,}6334
\end{bmatrix}
\]

\[
\lambda_{\max}=3{,}0404,\qquad
CI=0{,}0202,\qquad
CR=0{,}0348
\]

Zatem macierz subkryteriów pojemności dysku jest spójna.

\subsection*{4. Subkryteria dla kryterium Czas pracy na baterii}

Dla czasu pracy na baterii przyjęto, że dłuższy czas pracy jest korzystniejszy:
\[
12\text{ h} > 10\text{ h} > 8\text{ h} > 6\text{ h}
\]

\begin{table}[H]
\centering
\caption{Tabela subkryteriów dla kryterium Czas pracy na baterii}
\begin{tabular}{lcccc}
\toprule
 & \textbf{6 h} & \textbf{8 h} & \textbf{10 h} & \textbf{12 h} \\
\midrule
\textbf{6 h} & 1 & 1/3 & 1/5 & 1/7 \\
\textbf{8 h} & 3 & 1 & 1/3 & 1/5 \\
\textbf{10 h} & 5 & 3 & 1 & 1/3 \\
\textbf{12 h} & 7 & 5 & 3 & 1 \\
\bottomrule
\end{tabular}
\end{table}

\[
A_{S_4}=
\begin{bmatrix}
1 & 1/3 & 1/5 & 1/7 \\
3 & 1 & 1/3 & 1/5 \\
5 & 3 & 1 & 1/3 \\
7 & 5 & 3 & 1
\end{bmatrix}
\]

Suma kolumn:
\[
\sum S_{4,1}=16,\qquad
\sum S_{4,2}=9{,}3333,\qquad
\sum S_{4,3}=4{,}5333,\qquad
\sum S_{4,4}=1{,}6762
\]

Macierz znormalizowana:
\[
B_{S_4}=
\begin{bmatrix}
0{,}0625 & 0{,}0357 & 0{,}0441 & 0{,}0852 \\
0{,}1875 & 0{,}1071 & 0{,}0735 & 0{,}1193 \\
0{,}3125 & 0{,}3214 & 0{,}2206 & 0{,}1989 \\
0{,}4375 & 0{,}5357 & 0{,}6618 & 0{,}5966
\end{bmatrix}
\]

Wektor wag:
\[
W=
\begin{bmatrix}
0{,}0569\\
0{,}1219\\
0{,}2633\\
0{,}5579
\end{bmatrix}
\]

\[
A_{S_4}\cdot W=
\begin{bmatrix}
0{,}2299\\
0{,}4919\\
1{,}0994\\
2{,}3555
\end{bmatrix}
\]

\[
V=
\frac{A_{S_4}\cdot W}{W}
=
\begin{bmatrix}
4{,}0408\\
4{,}0362\\
4{,}1747\\
4{,}2222
\end{bmatrix}
\]

\[
\lambda_{\max}=4{,}1185,\qquad
CI=0{,}0395,\qquad
CR=0{,}0439
\]

Zatem macierz subkryteriów czasu pracy na baterii jest spójna.

\section{Przypisanie alternatyw do subkryteriów}

Każda alternatywa jest przypisana do odpowiedniego subkryterium:

\begin{table}[H]
\centering
\caption{Przypisanie tabletów do subkryteriów}
\begin{tabular}{lcccc}
\toprule
\textbf{Tablet} & \textbf{Cena} & \textbf{RAM} & \textbf{Dysk} & \textbf{Bateria} \\
\midrule
Tablet 1 & 2500 zł & 8 GB & 256 GB & 6 h \\
Tablet 2 & 3200 zł & 16 GB & 512 GB & 8 h \\
Tablet 3 & 4200 zł & 16 GB & 1000 GB & 10 h \\
Tablet 4 & 5200 zł & 32 GB & 1000 GB & 12 h \\
\bottomrule
\end{tabular}
\end{table}

\section{Ocena globalna alternatyw}

Ocenę globalną każdej alternatywy wyznaczono jako sumę iloczynów wag kryteriów oraz wag odpowiadających im subkryteriów:
\[
O(A_i)=\sum_{j=1}^{4} w_j \cdot s_{ij}
\]
gdzie:
\begin{itemize}
    \item \(w_j\) -- waga \(j\)-tego kryterium,
    \item \(s_{ij}\) -- waga subkryterium przypisanego \(i\)-tej alternatywie w obrębie \(j\)-tego kryterium.
\end{itemize}

Dla poszczególnych tabletów otrzymano:

\begin{table}[H]
\centering
\caption{Ocena globalna tabletów}
\begin{tabular}{lccccc}
\toprule
\textbf{Tablet} & \textbf{Cena} & \textbf{RAM} & \textbf{Dysk} & \textbf{Bateria} & \textbf{Ocena globalna} \\
\midrule
Tablet 1 & \(0{,}5579\cdot0{,}5579\) & \(0{,}2633\cdot0{,}1061\) & \(0{,}1219\cdot0{,}1061\) & \(0{,}0569\cdot0{,}0569\) & \(0{,}3554\) \\
Tablet 2 & \(0{,}5579\cdot0{,}2633\) & \(0{,}2633\cdot0{,}2605\) & \(0{,}1219\cdot0{,}2605\) & \(0{,}0569\cdot0{,}1219\) & \(0{,}2542\) \\
Tablet 3 & \(0{,}5579\cdot0{,}1219\) & \(0{,}2633\cdot0{,}2605\) & \(0{,}1219\cdot0{,}6334\) & \(0{,}0569\cdot0{,}2633\) & \(0{,}2288\) \\
Tablet 4 & \(0{,}5579\cdot0{,}0569\) & \(0{,}2633\cdot0{,}6334\) & \(0{,}1219\cdot0{,}6334\) & \(0{,}0569\cdot0{,}5579\) & \(0{,}3075\) \\
\bottomrule
\end{tabular}
\end{table}

Po wykonaniu działań liczbowych:
\[
O(A_1)=0{,}3554
\]
\[
O(A_2)=0{,}2542
\]
\[
O(A_3)=0{,}2288
\]
\[
O(A_4)=0{,}3075
\]

\section{Ranking alternatyw}

Na podstawie ocen globalnych otrzymano następujący ranking:
\begin{enumerate}
    \item Tablet 1 -- \(0{,}3554\)
    \item Tablet 4 -- \(0{,}3075\)
    \item Tablet 2 -- \(0{,}2542\)
    \item Tablet 3 -- \(0{,}2288\)
\end{enumerate}

Najwyższą ocenę globalną uzyskał \textbf{Tablet 1}. Oznacza to, że przy przyjętych preferencjach decydenta jest to najlepsza alternatywa.

\section{Wynik końcowy}

Zgodnie z poleceniem wykonano pełne obliczenia dla:
\begin{itemize}
    \item kryteriów najwyższego poziomu,
    \item subkryteriów każdego kryterium,
    \item oceny globalnej wszystkich alternatyw.
\end{itemize}

Na podstawie otrzymanych wag można stwierdzić, że:
\begin{itemize}
    \item najważniejszym kryterium jest \textbf{Cena},
    \item kolejne miejsca zajmują \textbf{Pamięć RAM}, \textbf{Pojemność dysku} oraz \textbf{Czas pracy na baterii},
    \item wszystkie analizowane macierze spełniają warunek zgodności \(CR<0{,}1\),
    \item najlepszą ocenę końcową uzyskał \textbf{Tablet 1}.
\end{itemize}

\section{Wnioski}

Metoda AHP pozwoliła uporządkować problem wyboru tableta i wyznaczyć ważność kryteriów decyzyjnych. Największą wagę otrzymało kryterium \textbf{Cena}, następnie \textbf{Pamięć RAM}, później \textbf{Pojemność dysku}, a najmniejszą \textbf{Czas pracy na baterii}.

Dla kryteriów najwyższego poziomu oraz dla wszystkich analizowanych grup subkryteriów otrzymano:
\[
CR<0{,}1
\]
co oznacza, że macierze porównań są spójne.

Na tej podstawie można uznać, że przedstawiony model decyzyjny jest poprawny, a najlepszą alternatywą przy przyjętych założeniach jest \textbf{Tablet 1}.

\end{document}